%% kapa: exam=1 round=24 year=2013 period=1 form=A subject=civil_law subject_ko=민법 seq=5 number=5
%% source: Q-Net 감정평가사 기출문제 (cst00309) artlSeq=5073802 file=extracted/1008/5073802/2013년 제24회 감정평가사 1차시험 기출문제/1교시(A형).hwp
%% conversion: kordoc 4.24.51 -> tools/gen1.py (auto); review=proofread-sonnet; flags=none
\begin{question}{5}
법인 아닌 사단에 관한 설명으로 옳지 않은 것은? (다툼이 있으면 판례에 의함)
\begin{choices}
  \item 교회의 대표자가 권한 없이 교회재산을 처분한 경우 권한을 넘은 표현대리에 관한 민법 제126조의 규정이 유추적용된다.
  \item 어떤 임야가 임야조사령에 의하여 동(洞)ㆍ리(里)의 명의로 사정(査定)된 경우 특별한 사정이 없는 한 그 동ㆍ리는 법인 아닌 사단이다.
  \item 종중은 자연발생적 단체로서 특별한 조직행위가 없더라도 법인 아닌 사단이 될 수 있다.
  \item 법인 아닌 사단이 당사자인 사건에서 대표자에게 적법한 대표권이 있는지의 여부는 법원의 직권조사사항이다.
  \item 법인 아닌 사단에는 사단법인에 관한 민법규정 중 법인격을 전제로 하지 않는 규정이 유추적용된다.
\end{choices}
\end{question}
\answer{1}
